\documentclass[11pt]{article}

\usepackage[margin=1in]{geometry}
\usepackage[T1]{fontenc}
\usepackage[utf8]{inputenc}
\usepackage{hyperref}
\usepackage{booktabs}
\usepackage{longtable}
\usepackage{array}
\usepackage{xcolor}
\usepackage{pdflscape}
\usepackage{xparse}

% Classifier key: monospace with safe break opportunities at slashes and
% underscores. We use the verbatim argument form (xparse 'v') so that the
% argument is read with catcodes that make \nolinkurl's break logic kick in
% at every / and _ regardless of how the source line was typed.
\NewDocumentCommand{\key}{v}{\nolinkurl{#1}}

\hypersetup{
  colorlinks=true,
  linkcolor=black,
  urlcolor=blue,
  pdftitle={metaDEBASS: Per-alert classifier outputs --- a coordination ask for collaborators},
}

\newcommand{\code}[1]{\texttt{#1}}
\newcommand{\filepath}[1]{\path{#1}}
\newcolumntype{L}[1]{>{\raggedright\arraybackslash}p{#1}}
\urlstyle{tt}
\setlength{\emergencystretch}{2em}

\title{metaDEBASS: Per-alert classifier outputs\\
\large A coordination ask for broker collaborators}
\author{Xianzhe (TZ) Tang \\ \texttt{tztang@bu.edu} \\ Boston University}
\date{April 27, 2026}

\begin{document}
\maketitle

\begin{abstract}
\noindent
metaDEBASS is an early-epoch meta-classifier for
Rubin/LSST. It combines broker classifier outputs with locally re-run
light-curve experts so we can rank transients for spectroscopic
follow-up after only a few detections. The strict rule is simple: a row
for $N$ detections may use only information available at those
$N$ detections, not anything learned later.

For several brokers we can query only the current object-level score.
That is useful for live triage, but not for training early-epoch models:
the current score may have been computed after many more detections than
the training row is allowed to see. This memo lists which classifier
outputs are safe now, which are blocked by timing, and what we would ask
broker teams to expose.
\end{abstract}

\section{Why per-alert outputs matter}

For one transient at one alert epoch, we want to predict a follow-up
priority: likely Type~Ia, interesting non-Ia, or neither. The training
table is keyed by \code{(diaObjectId, n\_det)}: one row per object per
detection count from the first detection to the latest.

Every column in that row has to be \emph{causally safe}. It must be
derivable from the alert stream and broker queries truncated at
$N$~detections. Light-curve features satisfy this by construction
because we re-extract them from truncated photometry. External
classifier scores satisfy it only when the broker exposes one of:

\begin{enumerate}
  \item a \emph{per-alert} endpoint that returns the score assigned to
    alert~$k$ when alert~$k$ was published, for $k=1,\dots,N$; or
  \item a \emph{timestamped} score history that lets us keep only scores
    whose timestamp is on or before detection~$N$.
\end{enumerate}

If the only available value is ``the broker's latest score for this
object today'', the score may reflect detections beyond~$N$. Using it at
$N=3$ when the broker had already seen $N=15$ is decision-time leakage.
The label is not leaking directly, but the model sees a future view of
the object, and that biases every metric we report.

\paragraph{Consequence for early-epoch prioritisation.}
Without per-alert outputs, we either drop the broker score from training
or re-run a local copy of the same method at the proper epoch cutoffs.
Local re-runs work for public models such as ALeRCE LC or SuperNNova, but
they duplicate broker infrastructure and do not reproduce a broker's
production preprocessing. They also cannot stand in for proprietary or
unreleased classifiers.

\section{Temporal tags used internally}

Every event written to the broker-event table carries a
\code{temporal\_exactness} tag with one of four values:

\begin{itemize}
  \item \textbf{exact\_alert}: the score was issued for a specific
    alert, so we can place it at the correct epoch.
  \item \textbf{rerun\_exact}: we re-ran the classifier locally with
    the photometry truncated to exactly $N$ detections.
  \item \textbf{static\_safe}: the output does not depend on later
    light-curve history, for example an image stamp classifier or a
    host-context flag.
  \item \textbf{latest\_object\_unsafe}: only a current object snapshot
    is available, with no per-alert or timestamped history. The
    no-leakage gold builder filters these rows by default
    (see \filepath{src/debass\_meta/ingest/gold.py}, around
    \code{select\_events\_asof()}).
\end{itemize}

The last bucket is the one we want to shrink.

\section{Reality check before the table}

I checked the names in the table against the local registry
(\filepath{src/debass\_meta/projectors/base.py}), local metrics, and
public documentation or papers. The classifier families are real; the
important distinction is whether \emph{we} currently use them in a
no-leakage training run.

External checks used for the table:
\begin{itemize}
  \item Fink ZTF fields \code{rf\_snia\_vs\_nonia},
    \code{snn\_snia\_vs\_nonia}, and \code{slsn\_score} are documented
    Fink science-module outputs:
    \url{https://doc.ztf.fink-broker.org/en/latest/broker/science_modules/}.
    Fink's LSST-facing CATS and early-SN classifiers are described in
    Fraga et al. 2024: \url{https://arxiv.org/abs/2404.08798}.
  \item ALeRCE exposes object probabilities and classifier metadata via
    its Python client, and its public papers describe the stamp, light
    curve, BHRF, and ATAT classifiers:
    \url{https://alerce.readthedocs.io/en/latest/apis.html} and
    \url{https://science.alerce.online/publications/}.
  \item Lasair Sherlock is a documented sky-context classifier:
    \url{https://lasair.readthedocs.io/en/main/core_functions/sherlock.html}.
  \item Pitt-Google documents its SuperNNova classifier and BigQuery
    output path:
    \url{https://pitt-broker.readthedocs.io/en/latest/broker/broker-overview.html}.
  \item Babamul/BOOM, ANTARES Superphot+, AMPEL SNGuess/FollowMe/FinalBet,
    ParSNIP, SuperNNova, ORACLE, SALT3, UPSILoN, and the
    \code{light-curve} package were checked against their papers or
    package documentation. For AMPEL specifically, Nordin et al. 2025
    defines three LSST workflows: SNGuess for infant transient filtering,
    FollowMe for follow-up sample selection, and FinalBet for final
    classification with priors:
    \url{https://arxiv.org/abs/2501.16511}. The original ZTF SNGuess
    paper is Miranda et al. 2022, A\&A 665, A99:
    \url{https://www.aanda.org/articles/aa/abs/2022/09/aa43668-22/aa43668-22.html}.
\end{itemize}

The AUC values below come from the latest per-expert metrics artifact in
this checkout, \filepath{reports/metrics/expert\_trust\_metrics\_safe\_v6e2.json},
unless a row says ``planned'' or ``registered''. That file has 11
measured per-expert reliability models. A v7 job script exists and is
set up to add \key{fink/slsn} and \key{ampel/snguess}, but this checkout
does not contain \filepath{reports/metrics/expert\_trust\_metrics\_safe\_v7.json}.
The local SNGuess model files are present under
\filepath{artifacts/local\_experts/snguess/}, but \code{xgboost} is not
installed in the local Python environments I checked, so the vendored
SNGuess wrapper does not load here.

\section{Current state, classifier by classifier}

Twenty-eight classifier or context keys are registered in
\filepath{src/debass\_meta/projectors/base.py}. The table lists each key,
its timing status, and what it is for. ``Used a lot'' here means one of
two things: either the broker produced many rows in the local
\filepath{data/silver/broker\_events.parquet}, or the latest safe metrics
file contains a measured per-expert reliability model with a nontrivial
test set.

\begin{itemize}
  \item \textbf{Survey}: ``ZTF'', ``LSST'', or ``any''.
  \item \textbf{Temporal}: one of \code{exact\_alert},
    \code{rerun\_exact}, \code{static\_safe}, or
    \code{latest\_unsafe}.
  \item \textbf{Per-alert?}: ``Yes'' if the broker exposes
    alert-resolved scores; ``Local'' if the score comes from our re-run
    on truncated photometry; ``N/A'' if the output is static context;
    ``No'' if only a latest object snapshot is available.
  \item \textbf{Current use}: measured reliability model, blocked by
    timing, feature-only, registered, or planned.
\end{itemize}

\begin{landscape}
\scriptsize
\setlength{\tabcolsep}{3pt}
\renewcommand{\arraystretch}{1.18}
\setlength{\emergencystretch}{2em}
\begin{longtable}{L{1.7cm} L{5.7cm} L{1.15cm} L{1.85cm} L{1.3cm} L{2.45cm} L{5.35cm}}
\toprule
\textbf{Broker} & \textbf{Classifier key} & \textbf{Survey} & \textbf{Temporal} & \textbf{Per-alert?} & \textbf{Current use} & \textbf{Notes} \\
\midrule
\endfirsthead
\toprule
\textbf{Broker} & \textbf{Classifier key} & \textbf{Survey} & \textbf{Temporal} & \textbf{Per-alert?} & \textbf{Current use} & \textbf{Notes} \\
\midrule
\endhead

% --- Fink ZTF ---
Fink & \key{fink/snn} & ZTF & exact\_alert & Yes & Reliability model (AUC 0.92) & Real Fink ZTF field \code{snn\_snia\_vs\_nonia}. SuperNNova Ia-vs-non-Ia score, one of our main ZTF signals: 509{,}532 silver rows across 8{,}602 objects; 23{,}556 safe test rows in v6e.2.\\
Fink & \key{fink/rf_ia} & ZTF & exact\_alert & Yes & Reliability model (AUC 0.81) & Real Fink ZTF field \code{rf\_snia\_vs\_nonia}. Random-forest early/rising SN~Ia score. Used often, but weaker than SNN: 254{,}766 silver rows; 23{,}556 safe test rows.\\
Fink & \key{fink/slsn} & ZTF & exact\_alert & Yes & Planned / registered & Real Fink ZTF field \code{slsn\_score}. Meant as an SLSN/SN-filter signal, not an Ia classifier. Newly registered here; absent from the local silver table and no local v7 metrics file is present. Fink returns $-1$ when there are not enough points or the object is too young.\\
\midrule
% --- Fink LSST ---
Fink LSST & \key{fink_lsst/snn} & LSST & exact\_alert & Yes & Reliability model (AUC 0.84) & Real LSST-era Fink SuperNNova-style score. Heavily backfilled in this checkout: 857{,}832 silver rows across 3{,}997 objects; 9{,}800 safe test rows. Useful, but the LSST label mix is still weak-label heavy.\\
Fink LSST & \key{fink_lsst/cats} & LSST & exact\_alert & Yes & Reliability model (AUC 0.85) & CATS is Fink's broad LSST classifier. Used heavily in silver (853{,}760 rows). In our ternary projection it mostly works as a broad transient/contaminant signal rather than a clean Ia probability.\\
Fink LSST & \key{fink_lsst/early_snia} & LSST & exact\_alert & Yes & Reliability model (AUC 0.86) & Real Fink early-SN~Ia score. Sparse by design: 12{,}895 silver rows and only 142 safe test rows, because it often does not trigger in the first few detections.\\
\midrule
% --- ALeRCE ZTF (LC) ---
ALeRCE & \key{alerce/lc_classifier_transient} & ZTF & latest\_unsafe & \textbf{No} & Blocked by timing & Real ALeRCE legacy light-curve classifier. We have 15{,}536 local silver rows, but they are current object snapshots, so the no-leakage build drops them. Use the local \key{alerce_lc} re-run instead.\\
ALeRCE & \key{alerce/lc_classifier_BHRF_forced_phot_transient} & ZTF & latest\_unsafe & \textbf{No} & Blocked by timing & Real ALeRCE BHRF forced-photometry transient classifier. We have 12{,}946 silver rows across 2{,}132 objects, but no per-alert history, so it is not used for safe training.\\
ALeRCE & \key{alerce/lc_classifier_BHRF_forced_phot_top} & ZTF & latest\_unsafe & \textbf{No} & Blocked by timing & Real ALeRCE BHRF top-level class grouping. Current snapshot only; 6{,}446 silver rows. Useful for live inspection, not safe as an early-epoch training feature.\\
ALeRCE & \key{alerce/LC_classifier_ATAT_forced_phot(beta)} & ZTF & latest\_unsafe & \textbf{No} & Blocked by timing & Real ALeRCE ATAT transformer classifier. It is common in the raw pull (119{,}396 silver rows), but the public query gives latest object probabilities, so we exclude it from no-leakage training by default.\\
% --- ALeRCE stamp ---
ALeRCE & \key{alerce/stamp_classifier} & ZTF & static\_safe & N/A & Registered / use carefully & Real ALeRCE image-stamp CNN. It is safe only when treated as first-alert image context. This local silver has mixed legacy tags: 9{,}475 static-safe rows plus 38{,}835 older latest-snapshot rows, so normalization matters before using it.\\
ALeRCE & \key{alerce/stamp_classifier_2025_beta} & ZTF & static\_safe & N/A & Registered / use carefully & Real newer ALeRCE stamp classifier. Same caution as the legacy stamp row: 3{,}384 static-safe rows plus 15{,}330 latest-snapshot rows in the local silver table.\\
ALeRCE & \key{alerce/stamp_classifier_rubin_beta} & LSST & static\_safe & N/A & Registered & Real Rubin/LSST stamp beta key from ALeRCE. Only 5 local silver rows for one object here, so it is not yet a measured contributor.\\
\midrule
% --- Lasair ---
Lasair & \key{lasair/sherlock} & any & static\_safe & N/A & Context feature & Real Lasair Sherlock sky-context output. It is not a posterior over transient classes; it gives host/star/AGN-like context. Local silver has 50 rows, so it is a useful feature when present but not a standalone classifier in our current run.\\
\midrule
% --- Babamul ---
Babamul & \key{babamul} & any & static\_safe & N/A & Context feature / registered & Real Babamul/BOOM broker path. We ingest six flags such as star, rock, stationary, and cross-survey match. This is contaminant context, not a probability classifier; no local measured reliability model yet.\\
\midrule
% --- Pitt-Google ---
Pitt-Google & \key{pittgoogle/supernnova_lsst} & LSST & exact\_alert & Yes & Reliability model (AUC 0.91) & Real Pitt-Google SuperNNova BigQuery output with \code{kafkaPublishTimestamp}. Used a lot in silver (866{,}838 rows across 3{,}998 objects); 6{,}552 safe test rows in v6e.2.\\
Pitt-Google & \key{pittgoogle/supernnova_ztf} & ZTF & latest\_unsafe & \textbf{No} & Blocked by timing & Real Pitt-Google ZTF SuperNNova table. Local silver has 22{,}028 rows, but the table path used here lacks an alert-time column, so we cannot align scores to $N$ detections safely.\\
Pitt-Google & \key{pittgoogle/upsilon_lsst} & LSST & static\_safe & N/A & Registered & UPSILoN is a real periodic-variable classifier. It is useful for rejecting variable-star contaminants, not for SN subtype probabilities. The Pitt-Google LSST table path is registered here, but this checkout has 0 local rows.\\
\midrule
% --- ANTARES ---
ANTARES & \key{antares/oracle} & any & exact\_alert & In principle & Dormant & ORACLE is a real 19-class LSST photometric classifier. ANTARES integration is registered, but no local rows or metrics are present here. Treat this as an integration target, not a current result.\\
ANTARES & \key{antares/superphot_plus} & any & exact\_alert & In principle & Reliability model in v6e.2 (AUC 0.66) & Superphot+ is real and runs via ANTARES. The safe v6e.2 metrics file has 180 test rows, but the current local \code{broker\_events.parquet} has no ANTARES rows; availability depends on the ANTARES backfill.\\
\midrule
% --- AMPEL ---
AMPEL & \key{ampel/snguess} & any & rerun\_exact & Local or API & Registered; not active here & SNGuess is real: Miranda et al. 2022 built it as a boosted-tree filter for young extragalactic transients, and AMPEL publishes ZTF candidates to TNS as AMPEL\_ZTF\_NEW. In this repo the same key can mean a local vendored T2BrightSNProb re-run or AMPEL T2 output. Locally verified status: 0 rows in silver/gold, and the wrapper does not load without \code{xgboost}. Do not cite the slide-deck AUC 0.95 until the v7 metrics file is present.\\
AMPEL & \key{ampel/parsnip_followme} & any & exact\_alert & Token API & Dormant & FollowMe is real, but it is not FinalBet. In Nordin et al. 2025, FollowMe uses the SNGuess-selected stream plus ParSNIP/LightGBM to choose a follow-up sample, mainly for uncertain SN~Ia vs SN~Ibc cases. The host/redshift/rate priors belong to FinalBet, which is not registered in this repo. No local AMPEL rows or safe metrics are present.\\
\midrule
% --- Local re-runs ---
Local & \key{parsnip} & any & rerun\_exact & Local & Dormant & ParSNIP is a real published transient light-curve model. Our wrapper needs compatible \code{model.pt} and \code{classifier.pkl}; no measured local contribution in this checkout.\\
Local & \key{supernnova} & any & rerun\_exact & Local & Reliability model (AUC 0.91) & Local SuperNNova re-run on truncated light curves. Used in the safe v6e.2 run with 34{,}793 test rows; also available on every row of the local DP1 50k snapshot file.\\
Local & \key{alerce_lc} & any & rerun\_exact & Local & Reliability model (AUC 0.88) & Local re-run of the ALeRCE light-curve classifier at proper epoch cutoffs. This is the safe substitute for the blocked ALeRCE object-snapshot probabilities; 30{,}887 safe test rows in v6e.2.\\
Local & \key{salt3_chi2} & any & rerun\_exact & Local & Reliability model (AUC 0.90) & SALT3 \(\chi^{2}\) Type-Ia consistency score using \code{sncosmo}. Physics-motivated and useful when the light curve has enough points; 37{,}675 safe test rows in v6e.2.\\
Local & \key{lc_features_bv} & any & rerun\_exact & Local & Reliability model (AUC 0.88) & Bazin/Villar parametric-feature expert using the \code{light-curve} package. Sparse on short DP1 light curves but measured on 39{,}998 safe test rows in v6e.2.\\
Local & \key{oracle_lsst} & LSST & rerun\_exact & Local & Dormant & Local ORACLE wrapper. The model is real, but our DP1 smoke test showed domain mismatch, with nearly all probability mass going to ``other''. Not used as a current result.\\

\bottomrule
\end{longtable}
\normalsize
\end{landscape}

\paragraph{Roll-up by current use.}
Measured in the latest safe metrics file: 11 keys
(\key{fink/snn}, \key{fink/rf_ia}, \key{fink_lsst/snn},
\key{fink_lsst/cats}, \key{fink_lsst/early_snia},
\key{pittgoogle/supernnova_lsst}, \key{antares/superphot_plus},
\key{supernnova}, \key{alerce_lc}, \key{salt3_chi2},
\key{lc_features_bv}). Blocked by timing: 6 rows, mainly ALeRCE
object snapshots plus Pitt-Google ZTF SuperNNova. Registered or dormant:
the rest. This is different from saying the classifiers are not real:
most are real external classifiers, but some are not usable in our
current no-leakage training build.

\section{Asks for the upcoming meeting}

The structural ask is the same for every broker team: \emph{do you
expose, or could you expose, the score assigned to alert~$k$ at the time
alert~$k$ was published?} Concretely:

\begin{description}
\item[ALeRCE.] Is there an internal per-alert table behind the
  light-curve classifiers that the public REST API does not currently
  expose? If not, would you be willing to publish a timestamped score
  history per object so downstream users can filter to ``score as of
  detection~$N$''? This is the largest gap for us. The local silver table
  has hundreds of thousands of ALeRCE object-probability rows, including
  ATAT and BHRF forced-photometry outputs, but the no-leakage build has
  to drop them.

\item[Pitt-Google ZTF table.] The LSST BigQuery table already carries
  \code{kafkaPublishTimestamp} per alert. Could the equivalent timestamp
  be added to the public ZTF SuperNNova table? Without it, we cannot
  reconstruct the per-alert score history even though the underlying
  model is a real and useful one.

\item[ANTARES and AMPEL.] We would like guidance on the preferred way to
  pull per-alert tags at scale. The classifiers are real; our current
  blocker is reliable bulk access and integration, not whether the
  methods exist. For AMPEL, the live Swagger page is reachable at
  \url{https://ampel.zeuthen.desy.de/api/live/docs}, and the OpenAPI
  schema exposes token-protected routes such as
  \code{/stock/ztf/\{ztf\_id\}} and \code{/stock/\{stock\_id\}/t2}.
  The local adapter still points at \code{/api/archive/object/.../t2},
  which returned 404 in my check. Before we claim AMPEL coverage, we
  should confirm the token, the correct route, and whether the returned
  T2 documents include enough state time or detection-count information
  to align scores to $N$ detections.

\item[Brokers running classifiers we already re-run locally.] A local
  substitute is useful, but it is not the broker's deployed score. It
  misses deployment-specific calibration, preprocessing, and versioning.
  If a broker can publish per-alert scores, we would rather use the
  broker's production output. ALeRCE LC is the immediate example.

\item[Brokers running proprietary or unreleased classifiers.] These are
  the cases where local substitution is impossible. Per-alert access
  matters most here.
\end{description}

\section{What we offer back}

In return, we can share, before any public release:

\begin{itemize}
  \item Reliability curves for each contributing classifier on a
    held-out matched cohort, computed under the no-leakage rule.
  \item Per-classifier ranking and calibration metrics
    (ROC--AUC, expected calibration error, Brier score) at
    $N=3,5,7,10$ detections.
  \item The labelled cohort itself, with truth separated by source
    quality: TNS spectroscopic, TNS-name-only, broker consensus, and
    host-context labels.
  \item The schema and naming contract our pipeline already uses, so a
    new per-alert endpoint can plug in without a redesign on either
    side.
\end{itemize}

\section{How to verify the numbers in this memo}

The registry can be checked locally in
\filepath{src/debass\_meta/projectors/base.py}. The tables used for the
row counts and timing checks are:

\begin{itemize}
  \item \filepath{data/silver/broker\_events.parquet}\par
    One row per broker--object--score, with \code{temporal\_exactness}.
  \item \filepath{data/gold/object\_epoch\_snapshots\_*.parquet}\par
    The no-leakage epoch tables. Per-classifier availability lives in
    the \code{avail\_\_<sanitized\_key>} columns.
  \item \filepath{reports/metrics/expert\_trust\_metrics\_safe\_v6e2.json}\par
    The latest per-expert safe metrics artifact present in this checkout.
\end{itemize}

The gating logic lives in
\filepath{src/debass_meta/ingest/gold.py} in
\code{select\_events\_asof()}. The flag
\code{--allow-unsafe-latest-snapshot} is for scoring-time experiments
only; it should not be used for training the per-expert reliability
models.

\section{Next steps}

After the meeting we should write down three things: which broker teams
to approach first, the exact endpoint or column we need from each one,
and who on our side will lead each conversation. A short project-state
README should travel with this memo so collaborators do not have to read
the codebase before they can help.

\end{document}
